Figure~\ref{fig:atlas-finite-time-blowup} shows the finite right endpoint of a maximal solution interval.

\begin{figure}[!htbp]
\centering
\begin{tikzpicture}[>=Stealth]
\begin{axis}[
  width=12.3cm,height=6.3cm,
  axis lines=left,
  ticklabel style={font=\small},label style={font=\small},
  xlabel={$t$},ylabel={$y(t)$},
  xmin=-1,xmax=1.25,ymin=0,ymax=10.5,
  xtick={-1,0,1},ytick={1,3,5,7,9}
]
\addplot[figblue,very thick,-{Stealth[length=2.4mm]},domain=-1:0.89,samples=180] {1/(1-x)};
\draw[figred,dashed,thick] (axis cs:1,0) -- (axis cs:1,9.8);
\node[figred,above left,fill=white,inner sep=1pt,font=\small]
  at (axis cs:0.98,9.6) {$T_+=1$};
\addplot[only marks,figblue,mark=*] coordinates {(0,1)};
\node[figblue,above left,font=\scriptsize] at (axis cs:0,1) {$y(0)=1$};
\end{axis}
\end{tikzpicture}
\caption{The solution $y(t)=1/(1-t)$ with initial value $y(0)=1$ satisfies $y^{\prime}=y^2$ and is defined on $(-\infty,1)$. The vertical dashed line marks the blow-up time, not part of the solution. The equation has smooth coefficients, but this trajectory cannot be continued through that time with a finite value.}
\label{fig:atlas-finite-time-blowup}
\end{figure}